\documentclass[a4paper,10pt]{article}
\usepackage[utf8]{inputenc} 
\usepackage[english]{babel}
\usepackage[scale=.7]{geometry}
\usepackage{graphicx}
\usepackage{amsmath,amssymb,amsthm}
\usepackage{hyperref}
\usepackage{polynomial}
\usepackage{enumitem}
\usepackage{todonotes}
\usepackage{tikz}
\usetikzlibrary{calc}

\numberwithin{equation}{section} 
\newtheorem{theorem}[subsubsection]{Theorem}
\newtheorem{corollary}[subsection]{Corollary}
\newtheorem{lemma}[subsection]{Lemma}
\newtheorem{proposition}[subsection]{Proposition}
\theoremstyle{definition}
\newtheorem{definition}[subsubsection]{Definition}
\newtheorem{remark}[subsection]{Remark}
\newtheorem{example}[subsection]{Example}

\title{Colouring minimal Cayley graphs: the road from Babai's conjectures to their resolution}

\author{Johan Voskamp}

\date{\today}
\begin{document}
\maketitle

\begin{abstract}
    In 1978 Babai conjectured that minimal Cayley graphs have a bounded chromatic number. In 1995 he changed his mind and conjectured the opposite.
\end{abstract}
\section{AI/Data usage}
    Used Claude to generate some TikZ's code with prompt:
    \textit{Give me standalone TikZ code for a Venn-style containment diagram: one large circle labeled 'no-lonely-colour graphs' fully containing a smaller circle labeled 'minimal Cayley graphs', 
    to show the proper-subset relationship. Use plain black outlines, light gray fill on the inner circle, and make sure the outer label doesn't overlap the inner circle. Compile it to check it renders before sending it to me.}

\begin{introduction}
\end{introduction}
\begin{preliminaries}
\end{preliminaries}
\begin{Babai's question, 19781995}
\end{Babai's question, 19781995}
\begin{Evidence in both directions}
\end{Evidence in both directions}
\begin{Necessary tools}
\end{Necessary tools}
\begin{No-lonely-colour graphs of large chromatic number}
\end{No-lonely-colour graphs of large chromatic number}
\begin{Minimal Cayley graphs of large chromatic number}
\end{Minimal Cayley graphs of large chromatic number}
\begin{Open questions}
\end{Open questions}
\begin{Appendix}
\end{Appendix}
\bibliography{Bibliography}
\bibliographystyle{plain}

\end{document}